\documentclass[letterpaper]{article}
\usepackage{amsmath,amssymb}
\usepackage{tikz-cd}
\usepackage[letterpaper,margin=1in]{geometry}
\setlength{\parindent}{0pt}
\newcommand{\noteheader}[1]{\par\medskip\noindent\textbf{#1}\ }

\begin{document}

\textbf{MATH 1564 --- Notes 09/29}

\bigskip

\textbf{Definition.} If $A$ is an $n\times n$ matrix in $\mathbb{R}^n$ such that
\[
Av=\lambda v,
\qquad v\ne0,
\qquad \lambda\in\mathbb{R},
\]
then $\lambda$ is an \textbf{eigenvalue} for $A$, and $v$ is an \textbf{eigenvector} for $A$.

\medskip

\noteheader{Question.} Is $\lambda$ real?
\[
\lambda=3,
\qquad
Av=\begin{bmatrix}9\\18\end{bmatrix}
=\lambda\begin{bmatrix}3\\6\end{bmatrix}.
\]

\bigskip

\noteheader{Theorem.} If $Av=\lambda v$, then $(\lambda^k v)$ solves
\[
X_{k+1}=AX_k.
\]

\noteheader{Proof.}
\[
X_{k+1}=\lambda^{k+1}v=A\lambda^kv
=\lambda^kAv=\lambda^k\lambda v.
\]

\clearpage

\noteheader{Theorem.} If $\{v_1,\ldots,v_n\}$ is a basis in $\mathbb{R}^n$ and
\[
Av_i=\lambda_i v_i,
\qquad i=1,2,\ldots,n,
\]
then
\[
\{(\lambda_1^kv_1),\ldots,(\lambda_n^kv_n)\}
\]
is a basis in the solution space of
\[
X_{k+1}=AX_k.
\]

\noteheader{Proof.} The solution space of $X_{k+1}=AX_k$ has dimension $n$, and each $(\lambda_i^kv_i)$ is a solution. To check independence, suppose
\[
c_1(\lambda_1^kv_1)+\cdots+c_n(\lambda_n^kv_n)=0.
\]
Set $k=0$. Then
\[
c_1v_1+\cdots+c_nv_n=0,
\]
so $c_1=\cdots=c_n=0$.

\clearpage

\noteheader{Representing a linear map by a matrix}

\medskip

Let $T:V\to W$ be a linear map where
\[
\dim V=n,
\qquad
\dim W=m.
\]
Let
\[
\beta=\{v_1,\ldots,v_n\}
\]
be a basis of $V$, and let
\[
\gamma=\{w_1,\ldots,w_m\}
\]
be a basis of $W$. Can we represent $T$ by a matrix?

Define coordinate maps
\[
S:V\to\mathbb{R}^n,
\qquad S(v)=[v]_\beta,
\]
and
\[
R:W\to\mathbb{R}^m,
\qquad R(w)=[w]_\gamma.
\]
The coordinate maps and $T$ fit into the diagram
\[
\begin{tikzcd}[column sep=large,row sep=large]
V \arrow[r,"S"] \arrow[d,"T"'] & \mathbb{R}^n \\
W \arrow[r,"R"'] & \mathbb{R}^m .
\end{tikzcd}
\]
Thus
\[
S(c_1v_1+\cdots+c_nv_n)=\begin{bmatrix}c_1\\\vdots\\c_n\end{bmatrix},
\]
and
\[
R(d_1w_1+\cdots+d_mw_m)=\begin{bmatrix}d_1\\\vdots\\d_m\end{bmatrix}.
\]

The map
\[
RTS^{-1}:\mathbb{R}^n\to\mathbb{R}^m
\]
is linear, so for some matrix $A$,
\[
RTS^{-1}(x)=Ax.
\]
The vector $Ae_i$ is the $i$th column of $A$. Since
\[
S^{-1}(e_i)=v_i,
\]
we have
\[
Ae_i=[T(v_i)]_\gamma.
\]

\clearpage

\noteheader{Theorem.}
\[
A=\begin{bmatrix}[T(v_1)]_\gamma&\cdots&[T(v_n)]_\gamma\end{bmatrix}.
\]

\bigskip

\noteheader{Example.} Let
\[
T:\mathbb{R}_{\leq2}[x]\to\mathbb{R}_{\leq1}[x],
\qquad T(f)=f',
\]
where $f'$ is the derivative of $f$. The usual basis in $\mathbb{R}_{\leq2}[x]$ is
\[
\beta=\{1,x,x^2\}.
\]
The coordinate map is
\[
S:\mathbb{R}_{\leq2}[x]\to\mathbb{R}^3,
\qquad
S(a_0+a_1x+a_2x^2)=\begin{bmatrix}a_0\\a_1\\a_2\end{bmatrix}.
\]
Then
\[
STS^{-1}=Ax,
\]
where
\begin{align*}
A
&=\begin{bmatrix}[T(1)]_\beta&[T(x)]_\beta&[T(x^2)]_\beta\end{bmatrix}\\
&=\begin{bmatrix}[0]_\beta&[1]_\beta&[2x]_\beta\end{bmatrix}
=\begin{bmatrix}0&1&0\\0&0&2\\0&0&0\end{bmatrix}.
\end{align*}

\bigskip

\noteheader{Question.} What is the smallest $k$ with $A^k=0$?
\[
k=3.
\]

\clearpage

\noteheader{Example.} Let $C$ be an $n\times n$ matrix. Define
\[
T:M_{n\times n}(\mathbb{R})\to M_{n\times n}(\mathbb{R}),
\qquad T(B)=CB.
\]
Is $T$ linear?
\begin{align*}
T(c_1B_1+c_2B_2)
&=C(c_1B_1+c_2B_2)\\
&=c_1CB_1+c_2CB_2\\
&=c_1T(B_1)+c_2T(B_2).
\end{align*}

\noteheader{How do we compute $\dim\ker T$ and $\dim\operatorname{Im}T$?}

\medskip

\noteheader{Subquestion exercise.} Let $n=2$ and
\[
C=\begin{bmatrix}a&b\\c&d\end{bmatrix}.
\]
A basis of $M_{2\times2}(\mathbb{R})$ is
\[
\left\{
\begin{bmatrix}1&0\\0&0\end{bmatrix},
\begin{bmatrix}0&1\\0&0\end{bmatrix},
\begin{bmatrix}0&0\\1&0\end{bmatrix},
\begin{bmatrix}0&0\\0&1\end{bmatrix}
\right\}.
\]
For this basis,
\begin{align*}
T\left(\begin{bmatrix}1&0\\0&0\end{bmatrix}\right)
&=\begin{bmatrix}a&0\\c&0\end{bmatrix}
=a\begin{bmatrix}1&0\\0&0\end{bmatrix}
+c\begin{bmatrix}0&0\\1&0\end{bmatrix},\\
\end{align*}

\clearpage

\begin{align*}
T\left(\begin{bmatrix}0&1\\0&0\end{bmatrix}\right)
&=\begin{bmatrix}0&a\\0&c\end{bmatrix},\\
T\left(\begin{bmatrix}0&0\\1&0\end{bmatrix}\right)
&=\begin{bmatrix}b&0\\d&0\end{bmatrix},\\
T\left(\begin{bmatrix}0&0\\0&1\end{bmatrix}\right)
&=\begin{bmatrix}0&b\\0&d\end{bmatrix}.
\end{align*}

With the coordinate map $S:M_{2\times2}(\mathbb{R})\to\mathbb{R}^4$,
\[
Ax=STS^{-1}(x),
\qquad
A=\begin{bmatrix}
a&0&b&0\\
0&a&0&b\\
c&0&d&0\\
0&c&0&d
\end{bmatrix}.
\]
Therefore,
\[
\dim\ker T=\dim\operatorname{Nul}A,
\qquad
\dim\operatorname{Im}T=\operatorname{rank}A.
\]

\clearpage

\noteheader{Subquestion.} If $C$ is invertible, is $T(B)=CB$ invertible?

\medskip

Yes.
\[
T^{-1}(B)=C^{-1}B,
\qquad
\ker T=\{0\}.
\]

\bigskip

\noteheader{Example.} Suppose $V$ has basis
\[
\beta=\{v_1,\ldots,v_n\}
\]
and another basis
\[
\gamma=\{u_1,\ldots,u_n\}.
\]
The coordinate-change map is represented by
\[
RS^{-1}(x)=Ax,
\qquad Ae_i=RS^{-1}(e_i).
\]
The change of coordinates is represented by
\[
\begin{tikzcd}[column sep=large,row sep=large]
V \arrow[r,"S"] \arrow[d,"\mathrm{id}"'] & \mathbb{R}^n \arrow[d,"RS^{-1}"] \\
V \arrow[r,"R"'] & \mathbb{R}^n .
\end{tikzcd}
\]

\clearpage

\noteheader{Example.} $\mathbb{R}^2$ has two bases
\[
\beta=\left\{
\begin{bmatrix}1\\-1\end{bmatrix},
\begin{bmatrix}1\\1\end{bmatrix}
\right\}
\]
and
\[
\gamma=\left\{
\begin{bmatrix}2\\1\end{bmatrix},
\begin{bmatrix}1\\0\end{bmatrix}
\right\}.
\]
The two coordinate maps are displayed as
\[
\begin{tikzcd}[column sep=large,row sep=large]
\mathbb{R}^2 \arrow[r,"S"] \arrow[d,"\mathrm{id}"'] & \mathbb{R}^2 \arrow[d,"RS^{-1}"] \\
\mathbb{R}^2 \arrow[r,"R"'] & \mathbb{R}^2 .
\end{tikzcd}
\]
Here
\[
S^{-1}(e_1)=v_1,
\qquad
S^{-1}(e_2)=v_2,
\qquad
S^{-1}=\begin{bmatrix}1&1\\-1&1\end{bmatrix},
\]
while
\[
R^{-1}=\begin{bmatrix}2&1\\1&0\end{bmatrix},
\qquad
R=\begin{bmatrix}0&1\\1&-2\end{bmatrix}.
\]
Therefore the coordinate-change matrix from $\beta$-coordinates to
$\gamma$-coordinates is
\[
RS^{-1}
=\begin{bmatrix}0&1\\1&-2\end{bmatrix}
\begin{bmatrix}1&1\\-1&1\end{bmatrix}
=\begin{bmatrix}-1&1\\3&-1\end{bmatrix}.
\]

\end{document}
